\documentclass[a4paper,10pt]{article}
\usepackage{paper-ru}
\usepackage{hyperref}

\def\thetitle{Аксиома единственности отражения}
\def\theauthors{Нина Лебедева и Антон Петрунин}

\hypersetup{colorlinks=true,
citecolor=black,
linkcolor=black,
anchorcolor=black,
filecolor=black,
menucolor=black,
urlcolor=black,
pdftitle={\thetitle},
pdfauthor={\theauthors}
}

%\overfullrule=100mm

\begin{document}

\title{\thetitle}
\author{\theauthors}
\date{}
\maketitle


\begin{abstract}
Приводится аксиоматика абсолютной плоскости, построенная на понятии метрического пространства с аксиомами существования геодезических, равенства треугольников по трём сторонам и единственности отражения.
\end{abstract}

От читателя потребуются базовые знания о метрических пространствах и аксиоматике абсолютной плоскости.

\begin{thm}{Теорема}\label{thm:main}
Любое нетривиальное трёхточечно однородное геодезическое пространство изометрично абсолютной плоскости,
если для любых двух различных его точек существует единственная нетождественная изометрия, фиксирующая обе точки.
\end{thm}

Начнём с расшифровки терминов.
Пространство, содержащее хотя бы две точки, называется \emph{нетривиальным}.

\emph{Геодезическая} в метрическом пространстве $M$ --- это изометричное отображение отрезка прямой в $M$.
Образ геодезической с концами $x$ и $y$ будет обозначаться $[xy]$; его также можно называть \emph{геодезической}.
Если любую пару точек $x,y\in M$ можно соединить геодезической, то $M$ называется \emph{геодезическим пространством}.
То есть требование геодезичности --- это ослабленный первый постулат Евклида: \textit{от всякой точки до всякой точки можно провести прямую линию}.

\emph{Изометрией} $M$ называется биективное отображение из $M$ в себя, сохраняющее расстояния.
На абсолютной плоскости единственная нетождественная изометрия, фиксирующая пару  различных точек, --- это, конечно, отражение в прямой, содержащей обе точки.

Трёхточечная однородность --- это по сути признак равенства треугольников по трём сторонам.
Метрическое пространство $M$ \emph{трёхточечно однородно}, если для любой шестёрки точек $x_1,x_2,x_3,y_1,y_2,y_3\in M$ равенства $|x_i-x_j|=|y_i-y_j|$ для всех $i$ и $j$ влекут существование изометрии $M$, отображающей $x_i$ в $y_i$ для каждого $i$.
Здесь $|p-q|$ обозначает расстояние между точками $p,q\in M$.
Точки в тройке не обязаны быть различными; в частности, трёхточечная однородность влечёт двуточечную и обычную (одноточечную) однородность.

Идея использовать в таком контексте признак равенства треугольников по трём сторонам встречается у Кароля Борсука и Ванды Шмелевой \cite[part I chap. II § 4]{borsuk-szmielew}, а также у Александра Даниловича Александрова \cite[гл. 2, § 8]{alexandrov};
в обоих случаях равенство углов определяется через равенство сторон вспомогательных треугольников.
Особенность нашего подхода в том, что исходным объектом служит геодезическое пространство, и все свойства плоскости выводятся из условий на его изометрии.
Похожая аксиоматика была предложена Джорджем Биркгофом \cite{birkhoff1932}, но кроме чисто метрических понятий он использовал измерения углов.
В нашем подходе естественно сплетаются геометрия метрических пространств и основания геометрии, и он может оказаться востребованным в преподавании.

Все идеи в доказательстве элементарны, и многие очень древние.
Заметим, что мы не предполагаем локальную компактность пространства, иначе наша теорема была бы немедленным следствием теоремы Жака Титса~\cite[IV.E § 3.59 cor. 2]{tits}, классифицирующей локально компактные двуточечно однородные геодезические пространства.
Толчком к написанию заметки послужил другой результат авторов~\cite{lebedeva-petrunin}.

\section*{Доказательство}

Пусть $M$ удовлетворяет условиям нашей теоремы.

\begin{wrapfigure}{r}{35mm}
\centering
\vskip-15mm
\includegraphics{mppics/pic-10}
\vskip-0mm
\end{wrapfigure}

\parbf{Пара расстояний.}
Пусть $p,q\in M$ --- две различные точки. Пара расстояний $r=|p-x|$ и $s=|q-x|$ определяет точку $x$ с точностью до двух возможностей.

Действительно, существует единственная нетождественная изометрия, скажем $\sigma$, фиксирующая $p$ и~$q$.
Если $|p-y|=r$ и $|q-y|=s$, то по трёхточечной однородности найдётся изометрия, фиксирующая $p$ и $q$ и переводящая $x$ в $y$.
То есть $y=x$, если изометрия тождественна, а иначе $y=\sigma(x)$ --- третьего не дано.

\parbf{Локальная единственность геодезических.}
Покажем, что существует $\eps_0>0$ такое, что любые две точки $x$ и $y$ на расстоянии меньше $\eps_0$ соединяются единственной геодезической.

\emph{Серединой} пары точек $x$ и $y$ будет называться точка $m$, для которой
\[|x-m|=|y-m|=\tfrac12\cdot|x-y|.\]
Если середина единственна, то любая геодезическая $[xy]$ её содержит.
Если же середины единственны для всех пар точек на расстоянии меньше $\eps_0$, то, итерируя переход к серединам, получаем единственность геодезических длины меньше $\eps_0$.

По наблюдению выше, у пары точек может быть не больше двух середин.
Предположим, что середина не всегда единственна, скажем у $p$ и $q$ нашлось две середины $m_1$ и $m_2$; положим $\eps_0=|m_1-m_2|$.

Пусть $0<|x-y|<\eps_0$.
Выберем точки $\bar x\z\in [m_1p]$ и $\bar y\in [m_1q]$ на расстоянии $\tfrac12\cdot|x-y|$ от $m_1$.
По неравенству треугольника, $|\bar x-\bar y|=|x-y|$.

\begin{wrapfigure}{r}{35mm}
\centering
\vskip-0mm
\includegraphics{mppics/pic-20}
\vskip-0mm
\end{wrapfigure}

Предположим, что у $x$ и $y$ нашлось две середины, скажем $n_1$ и $n_2$.
По трёхточечной однородности найдётся изометрия
$\iota$, переводящая $x\mapsto \bar x$, $y\mapsto \bar y$ и $n_1\mapsto m_1$.
Тогда (опять же по неравенству треугольника) точка $m_3=\iota(n_2)$ была бы серединой $p$ и $q$.
Заметим, что $m_3\ne m_1$ и $m_3\ne m_2$; последнее верно, поскольку
\[|m_3-m_1|\le |x-y|<\eps_0=|m_1-m_2|.\]
То есть середины $m_1$, $m_2$ и $m_3$ различны, что невозможно.


\parbf{Отсутствие ветвлений.}
Покажем, что \textit{геодезические в $M$ не ветвятся}, то есть если
две геодезические $[xz]$ и  $[xz']$ содержат нетривиальную геодезическую $[xy]$, то $[xz]$ содержит $[xz']$ или наоборот $[xz]$ содержится в $[xz']$.


Если бы геодезические ветвились, то, укоротив их и приблизив $x$ к месту ветвления, можно было бы считать, что $|x-z|=|x-z'|<\eps_0$ и $z\ne z'$.

{

\begin{wrapfigure}{R}{35mm}
\centering
\vskip-0mm
\includegraphics{mppics/pic-30}
\vskip-0mm
\end{wrapfigure}

По трёхточечной однородности найдётся изометрия $\sigma$, фиксирующая $x,y$ и переводящая $z$ в $z'$.
Но есть ещё нетождественная изометрия, скажем $\tau$, фиксирующая $x$ и $z$; она обязана фиксировать все точки на единственной геодезической $[xz]$, включая~$y$.
Итак, $\sigma$ и $\tau$ --- две различные нетождественные изометрии, фиксирующие $x$ и $y$ --- противоречие.

}

\parbf{Прямые.}
Любую геодезическую можно продолжить до \emph{локальной геодезической} $\gamma\colon\mathbb R\to M$; это значит, что достаточно короткие участки $\gamma$ являются геодезическими.
Для этого совместим короткий конец данной геодезической с началом более длинной.
По локальной единственности они совпадут на общем участке, что даст продолжение.
Повторяя это с одним и тем же малым шагом в обе стороны, получим $\gamma$.
Отсутствие ветвлений обеспечивает единственность продолжения локальной геодезической по любому её нетривиальному участку.

Пусть $d$ --- диаметр $M$; докажем, что $d=\infty$.
Следующее рассуждение заимствовано у Гарретта Биркгофа~\cite{birkhoff1944}.

Заметим, что в $M$ найдётся геодезическая любой длины меньше $d$.
По двуточечной однородности пространства и отсутствию ветвлений эти геодезические можно совместить с локальной геодезической $\gamma$, построенной выше.
Значит
\[|\gamma(t)-\gamma(s)|=|t-s|\quad\text{при}\quad |t-s|<d.
\eqlbl{eq:|t-s|}\]

Пусть $d<\infty$, тогда по непрерывности \ref{eq:|t-s|} выполняется и при $|t-s|=d$.
Выберем малое $\eps>0$ и положим
\[p=\gamma(0),
\quad x=\gamma(d+\eps),
\quad r=|p-x|,
\quad y=\gamma(r)
\quad\text{и}\quad
z=\gamma(\eps).\]
По неравенству треугольника $d-\eps\le r\le d$.

\begin{wrapfigure}{r}{35mm}
\centering
\vskip-8mm
\includegraphics{mppics/pic-40}
\vskip-0mm
\end{wrapfigure}

Точки $x$ и $y$ равноудалены от $p$, поэтому существует изометрия $\sigma$, фиксирующая $p$ и переставляющая $x$ и $y$.
Она обращает геодезическую $[xy]$, а значит фиксирует её середину, скажем $m$.
По нашим условиям существует другая нетождественная изометрия, фиксирующая $p$ и $z$;
она обязана фиксировать единственную геодезическую $[pz]$, а значит и все точки на $\gamma$.
Получились две различные нетождественные изометрии, фиксирующие $p$ и $m$ --- противоречие.

Итак, $d=\infty$, и каждый участок локальной геодезической $\gamma$ является геодезической;
то есть $\gamma$ --- геодезическая; она задаётся изометрическим вложением всей вещественной прямой в $M$.
Такие геодезические будут называться \emph{прямыми}.

Проверим единственность геодезической между точками $x$ и $y$.
Предположим, таких нашлось две.
Продолжим одну из них за $y$ до точки $z$.
Другая вместе с тем же участком $[yz]$ также будет геодезической $[xz]$.
Получилось ветвление, что невозможно.

\begin{wrapfigure}{r}{35mm}
\centering
\vskip-5mm
\includegraphics{mppics/pic-50}
\vskip-0mm
\end{wrapfigure}

Итак, через две различные точки проходит единственная прямая.
Далее, существует единственная нетождественная изометрия, фиксирующая все точки на прямой; она будет называться \emph{отражением} относительно этой прямой.
Кроме того, из предположений теоремы следует, что отражение является инволюцией; то есть обратно самому себе.

\parbf{Центральная симметрия.}
Выберем точку $m$.
Для данной точки $x$ продолжим геодезическую $[xm]$ за $m$ на то же расстояние.
Обозначим противоположный конец через $x'$ и положим $m'=m$.
Поскольку прямая единственна, отображение $\theta\colon x\mapsto x'$ определено однозначно.
Покажем, что это изометрия.
Построенную изометрию будем называть \emph{центральной симметрией} относительно $m$.

\begin{figure}[h!]
\centering
\vskip-0mm
\includegraphics{mppics/pic-60}
\vskip-0mm
\end{figure}

Из равенства $x''=x$ следует биективность $\theta\colon x\mapsto x'$, так что достаточно проверить равенство
\[|x'-y'|=|x-y|\eqlbl{eq:xy=xy}\]
для любых двух точек $x$ и $y$.
При этом можно считать, что $m$, $x$ и $y$ попарно различны, иначе вовсе нечего доказывать.

Выберем $\eps>0$ меньше $|m-x|$ и $|m-y|$ и точки $\bar x\in[mx]$ и $\bar y\in[my]$ на расстоянии $\eps$ от $m$.
По трёхточечной однородности найдётся изометрия $\sigma$, фиксирующая $m$ и переставляющая $\bar x$ и $\bar y'$.
Изометрия переводит прямые в прямые, поэтому она переставляет $\bar x'$ и $\bar y$.
В частности, $|\bar x-\bar y|=|\bar x'-\bar y'|$.
Снова по трёхточечной однородности найдётся изометрия, фиксирующая $m$ и переводящая $\bar x\mapsto \bar x'$ и  $\bar y\mapsto \bar y'$.
По единственности продолжения она переводит $x\mapsto x'$ и $y\mapsto y'$;
отсюда \ref{eq:xy=xy}.

\begin{wrapfigure}{r}{35mm}
\centering
\vskip-8mm
\includegraphics{mppics/pic-100}
\vskip-0mm
\end{wrapfigure}

\parbf{Отражения.}
Пусть $\sigma$ --- отражение относительно прямой $\ell$.
Покажем, что \textit{$\sigma$ фиксирует только точки на $\ell$}.

Выберем две различные точки $p$ и $q$ на $\ell$, и пусть $m$ --- их середина.
Предположим, что $\sigma(u)=u$ для некоторой точки $u\notin\ell$.
Тогда $\sigma$ фиксирует все точки на ломаной $[pu]\cup[uq]$.
По единственности прямых, ломаная не проходит через $m$.
При движении точки $x$ вдоль ломаной от $p$ к $q$ разность $|q-x|-|p-x|$ меняет знак.
Значит, на ломаной найдётся точка $v$, равноудалённая от $p$ и $q$.

Пусть $v'$ центрально симметрична точке $v$ относительно $m$.
Заметим, что $|p-v|=|q-v|=|p-v'|\z=|q-v'|$
и $v'\ne v$, ведь $v\ne m$.
По трёхточечной однородности существует изометрия, переводящая $v$ в $v'$ и фиксирующая $p$ и $q$.
По единственности, это должна быть $\sigma$;
то есть $\sigma(v)=v'\ne v$ --- противоречие.

\parbf{Серединный перпендикуляр.}
Как и раньше, $\sigma$ --- отражение относительно прямой $\ell$.
Пусть $v\notin\ell$ и $w=\sigma(v)$; мы уже знаем, что $v\ne w$.

\begin{wrapfigure}{r}{35mm}
\centering
\vskip-0mm
\includegraphics{mppics/pic-110}
\vskip-0mm
\end{wrapfigure}

Докажем, что
\[x\in \ell\quad\Leftrightarrow\quad|v-x|=|w-x|;\]
прямая импликация очевидна.


Пусть $m$ --- середина $v$ и $w$ и $\theta$ --- центральная симметрия относительно $m$.
Заметим, что $\sigma(m)\z=m$, и значит, $m\in\ell$.
Далее, $\rho=\theta\circ\sigma$ есть отражение относительно прямой $vw$;
действительно, $\rho$ фиксирует $v$ и $w$ и разворачивает~$\ell$.

Пусть $|v-x|=|w-x|$.
Если $x\notin \ell$, то $\sigma(x)\ne x$ и $\theta(x)\ne x$.
Все три точки $x$, $\sigma(x)$ и $\theta(x)$ лежат на одинаковых расстояниях от $v$ и $w$;
значит $\sigma(x)=\theta(x)$ и $x=\rho(x)$, то есть $x$ лежит на прямой $vw$.
Поскольку $|v-x|=|w-x|$, получаем $x=m\in\ell$ --- противоречие.

\parbf{Полуплоскости.}
Пусть
\begin{align*}
V&=\set{x\in M}{|v-x|<|w-x|}
&&\text{и}&
W&=\set{x\in M}{|v-x|>|w-x|}.
\end{align*}
Поскольку $V\ni v$ и $W\ni w$, оба множества непусты.
Отрезок, соединяющий точки разных классов, обязан пересекать $\ell$,
ведь разность расстояний до $v$ и $w$ меняет на нём знак.

Покажем, что  $V$ и $W$ \emph{выпуклы}, то есть вместе с любой парой точек содержат геодезическую с их концами.
Действительно, если $x\in V$, то по неравенству треугольника $[vx]\subset V$.

Пусть теперь $x,y\in V$.
Ломаная $[xv]\cup[vy]$ лежит в $V$ и потому не пересекает $\ell$.
На этой ломаной расстояния до $x$ и $\sigma(x)$ не могут сравняться; это возможно только на $\ell$.
В её начале $0=|x-x|<|\sigma(x)-x|$, поэтому и в конце $|x-y|<|\sigma(x)-y|$.
Если бы на $[xy]$ нашлась точка $z\in\ell$, то
\[|\sigma(x)-y|\le |\sigma(x)-z|+|z-y|=|x-z|+|z-y|=|x-y|\]
--- противоречие.
То есть $[xy]$ не пересекает $\ell$ и по непрерывности целиком лежит в $V$.
Выпуклость $W$ доказывается так же.

Заметим, что отсюда следует аксиома Паша: \textit{прямая, не проходящая через вершины треугольника и пересекающая одну из его сторон, обязана пересечь и одну из двух других сторон.}

\parbf{Концовка.}
Остаётся проверить, что из вышесказанного следует выполнение аксиоматики абсолютной плоскости;
эта проверка довольно простая, но она зависит от выбранной системы аксиом.

Аксиоматику абсолютной плоскости можно получить из аксиоматики Тарского~\cite{tarski}, исключив аксиому A8 (она эквивалентна пятому постулату Евклида);
перевод понятий на метрический язык прямолинеен.
Вместо схемы аксиом A13 можно взять более сильную полную аксиому непрерывности \cite[с. 18]{tarski}; она выполняется, поскольку прямые изометричны $\mathbb R$.
%Аксиома A12 следует из параграфа о серединном перпендикуляре.

\parit{Мы использовали ИИ при обсуждении доказательства, поиске литературы и редактировании текста.}

{\sloppy
\def\emph{\textit}
\printbibliography[heading=bibintoc]\fussy}
\end{document}

\newpage

\parbf{Неравенство для середин.}
Рассмотрим пару геодезических $[p_0q_0]$ и $[p_1q_1]$;
пусть $u_0$ и $u_1$ --- их середины.
Покажем, что
\[|u_0-u_1|\le \frac{|p_0-p_1|+|q_0-q_1|}2.\eqlbl{eq:mid}\]
Доказательство очень близко к доказательству Адриена Лежандра \cite{legendre1833} того, что сумма углов треугольника на абсолютной плоскости не превосходит $\pi$.

\begin{figure}[h!]
\centering
\vskip-0mm
\includegraphics{mppics/pic-70}
\vskip-0mm
\end{figure}

Можно считать, что $u_0\ne u_1$, иначе нечего доказывать.
Продолжим конфигурацию точек, применяя центральную симметрию,
полагая, что $u_{i+1}$, $p_{i+1}$ и $q_{i+1}$ центрально симметричны $u_{i-1}$, $q_{i-1}$ и $p_{i-1}$ относительно $u_i$.
По неравенству треугольника получаем
\[|u_0-u_n|\le|u_0-p_0|+|p_0-p_1|+\ldots+|p_{n-1}-p_n|+|p_n-u_n|.\]
Заметим, что $u_i$ есть середина $[p_iq_i]$ и отрезки, одинаково отмеченные на рисунке, имеют равную длину.
В частности, длины $|p_i-p_{i+1}|$ поочерёдно равны $|p_0-p_1|$ и $|q_0-q_1|$.
Если $n$ чётно, то справа в неравенстве стоит два равных слагаемых $|u_0-p_0|=|u_n-p_n|$,
и по $\tfrac n2$ слагаемых $|p_0-p_1|$ и $|q_0-q_1|$.
Кроме того, точки $u_0,\dots,u_n$ лежат на одной прямой, и поэтому $|u_0-u_n|=n\cdot |u_0-u_1|$.
Значит неравенство можно переписать как
\[n\cdot |u_0-u_1|\le 2\cdot |u_0-p_0|+\tfrac n2\cdot |p_0-p_1|+ \tfrac n2\cdot |q_0-q_1|.\]
Значение $|u_0-p_0|$ не зависит от $n$;
поскольку $n$ можно взять произвольно большим, получаем \ref{eq:mid}.

\parbf{Зависимость геодезической от концов.}
Зафиксируем точку $o$.
Для данной точки $x$ и $t\z\in[0,1]$
обозначим через $x_t$ точку на геодезической $[ox]$, такую, что $|o-x_t|\z=t\cdot|o-x|$.

Повторяя деление пополам и применяя \ref{eq:mid}, можно получить следующее неравенство:
\[|x_t-y_t|\le  t\cdot|x-y|.\]
Сначала оно проверяется для двоичных дробей $t$, а затем обобщается на все $t$ по непрерывности.


\begin{wrapfigure}[3]{r}{30mm}
\centering
\vskip-6mm
\includegraphics{mppics/pic-80}
\vskip-0mm
\end{wrapfigure}

Заметим, что
\begin{align*}
|x_t-y_s|
&\le|x_t-y_t|+|y_t-y_s|\le
\\
&\le t\cdot|x-y|+|t-s|\cdot |o-y|.
\end{align*}
Отсюда немедленно следует, что отображение $(x,t)\z\mapsto x_t$ непрерывно.%
\footnote{Ещё отсюда можно извлечь, что геодезическая $[xy]$ непрерывно (в смысле Хаусдорфа) зависит от концов.
Примеры геодезических пространств с единственными геодезическими, в которых геодезические не зависят непрерывно от концов, обсуждаются в книге Мартина Бридсона и Андре Хаефлигера \cite[I.3.14]{bridson-haefliger}.}

\parbf{Компактность веера.}
\emph{Веером} геодезической $[vw]$ относительно точки $o$ будет называться объединение геодезических $[ox]$ при $x\in[vw]$.
Заметим, что такой веер является образом непрерывного отображения $(x,t)\mapsto x_t$ при $x\in[vw]$ и $t\in [0,1]$,
и \textit{веер компактен}, поскольку множества $[vw]$ и $[0,1]$ компактны.

\parbf{Локальная компактность пространства.}
Напомним, что пространство локально компактно, если у каждой его точки есть замкнутая компактная окрестность.

\begin{figure}[ht!]
\centering
\vskip-0mm
\includegraphics{mppics/pic-90}
\vskip-0mm
\end{figure}

Выберем прямую $\ell$, проходящую через точку $o$.
Пусть $\sigma$ --- нетождественная изометрия, фиксирующая две различные точки $\ell$, а значит и все точки на прямой.
Выберем точку $y$, которую $\sigma$ перемещает; в частности, $y\notin\ell$.

Возьмём на $\ell$ точки $x$ и $z$ по разные стороны от $o$.
По единственности прямых, ломаная из звеньев $[xy]$ и $[yz]$ не проходит через $o$.
Поскольку ломаная компактна, расстояние от $o$ до неё положительно.
Выберем $\eps>0$ меньше этого расстояния.

Пусть $F$ --- объединение двух вееров над отрезками $[xy]$ и $[yz]$ относительно $o$;
как  объединение двух вееров оно компактно.
Покажем, что $F\cup\sigma(F)$ содержит окрестность $o$.


Пусть $0<|o-q|=r<\eps$.
По неравенству треугольника
\[|x-o|-r\le|x-q|\le|x-o|+r.\]
Для каждой точки $w$ на ломаной отметим точку $\bar w\in[ow]$ на расстоянии $r$ от~$o$.
Иначе говоря, $\bar w=w_t$, где $t=\tfrac r{|o-w|}$, точка $w_t$ определена чуть выше.
Из вышесказанного следует, что отображение $w\mapsto \bar w$ непрерывно на $[xy]\cup [yz]$.
При этом расстояние $|x-\bar w|$ непрерывно меняется от $|x-o|-r$ до $|x-o|+r$.
Значит, в некоторый момент $|x-\bar w|=|x-q|$.

Поскольку у точек $\bar w$ и $q$ одинаковая пара расстояний до $o$ и $x$,
$q=\bar w$ или $q=\sigma(\bar w)$, и в любом случае $q\in F\cup\sigma(F)$.

Итак, компактное множество $F\cup\sigma(F)$ содержит окрестность точки $o$ --- локальная компактность доказана.

\parbf{Концовка доказательства.}
Следующее утверждение является следствием теоремы Жака Титса~\cite{tits}; оно приведёт нас к финишу.

\parit{Любое нетривиальное локально компактное трёхточечно однородное геодезическое пространство изометрично евклидову пространству, сфере или пространству Лобачевского некоторой конечной размерности.}

Действительно, сфера исключается, поскольку диаметр $M$ бесконечен.
Размерность один также невозможна: изометрия прямой, фиксирующая две различные точки, тождественна.
При размерности больше двух существует бесконечно много изометрий, фиксирующих две заданные точки: можно вращать пространство вокруг проходящей через них прямой.
Поэтому размерность равна двум и наша теорема доказана.






\parbf{О другой концовке доказательства.}
Есть более элементарный способ закончить доказательство.
Можно свести задачу к проверке уже упомянутых аксиом Джорджа Биркгофа \cite{birkhoff1932}, адаптированных для абсолютной геометрии (то есть всех аксиом, кроме последней, в \cite{petrunin}).
Используя рассуждение из доказательства локальной компактности, можно построить транспортир,
то есть научиться измерять величину угла и доказать признак равенства треугольников по двум сторонам и углу между ними (это основная аксиома в системе Биркгофа).

Такая концовка более элементарна; она обходится без теоремы классификации, но несколько длиннее.
В любом случае концовка доказательства зависит от того, какое определение абсолютной плоскости вы выбрали.


В дальнейшем развитии метрического подхода можно выделить два направления.
В одном из них, начатом Гарреттом Биркгофом~\cite{birkhoff1944}, усиливается требование однородности:
пространство называется \emph{полностью однородным}, если любая изометрия между его подмножествами продолжается до изометрии всего пространства.
Подмножества здесь могут быть и бесконечными. В нашей недавней работе~\cite{lebedeva-petrunin} получена классификация таких геодезических пространств без других предположений. Помимо евклидовых пространств, сфер и пространств Лобачевского, в списке оказываются так называемые универсальные метрические деревья конечной валентности.

В другом направлении сохраняют трёхточечную однородность, но предполагают локальную компактность. Это позволяет применять теоремы о локально компактных топологических группах.
Мы сведём нашу теорему к одной из таких теорем --- теореме Жака Титса~\cite{tits}.


%Из неё также можно получить аксиоматику евклидовой геометрии, добавив в аксиомы, что пространства с обычной и удвоенной метрикой изометричны.

Отметим, что метрические системы аксиом используют вещественные числа как готовый блок, и это позволяет сделать их более экономными.
Первая подобная система аксиом была предложена Джорджем Биркгофом \cite{birkhoff1932} --- четыре аксиомы для евклидовой геометрии.
Если потребовать выполнения аксиомы подобия (последней аксиомы) только при коэффициенте $k=1$, то получатся аксиомы абсолютной геометрии.

В отличие от нашей системы, аксиоматика Биркгофа позволяет перейти к содержательной геометрии почти без предварительной подготовки.
Поэтому её часто используют в учебниках, включая учебник второго автора \cite{petrunin}.
Обе аксиоматики применимы только в двумерном случае.
